%! Graphing Quadratic Functions & Transformations — Unit 2, Lesson 2.1 — Student Packet Cover Sheet
% 12pt like every other student component. The banner is \coverbanner, which measures the title
% block and sizes the band to it. Three scored rows: Warm-Up, Guided Notes & Practice, Homework.
% Homework is due the first class after two study halls — course policy of record (2026-09-03) —
% never "next class". \namedateperiod appears HERE AND NOWHERE ELSE in the packet.
\documentclass[12pt]{article}
\usepackage{algebra2-article}
\usepackage{algebra2-boxes}

\begin{document}

\coverbanner{Unit 2 \quad Quadratic Functions}{Lesson 2.1 \quad Graphing Quadratic Functions \& Transformations}

\namedateperiod
\vspace{0.12in}

\boxguard[14]
\begin{learningtargetbox}
\small
\begin{enumerate}[leftmargin=1.5em, itemsep=3pt]
  \item \ldots{} describe any parabola as a \textbf{transformation} of the \textbf{parent
        function} $y=x^2$ --- \textbf{shift}, \textbf{stretch}, \textbf{reflection} --- knowing
        that \emph{subtracting inside the bracket moves the graph right}.
  \item \ldots{} read the \textbf{vertex} $(h,k)$ and the \textbf{axis of symmetry} $x=h$
        straight out of \textbf{vertex form} $a(x-h)^2+k$.
  \item \ldots{} read the $y$-intercept $(0,c)$ out of \textbf{standard form} $ax^2+bx+c$, find
        its axis with $x=-\frac{b}{2a}$, and substitute back to get the vertex.
  \item \ldots{} read the \textbf{zeros} out of \textbf{intercept (factored) form}
        $a(x-p)(x-q)$, place the axis at their midpoint, and show that all three forms of one
        curve agree.
\end{enumerate}
\end{learningtargetbox}

\vspace{0.10in}

\boxguard[14]
\begin{tocbox}
\small
\renewcommand{\arraystretch}{1.35}
\begin{tabularx}{\linewidth}{@{} c l X r @{}}
\rowcolor{forestbg}
\textbf{\#} & \textbf{Component} & \textbf{Description} & \textbf{Score} \\
\midrule
1 & Warm-Up & Expanding two expressions that turn out to be the same, moving the V of $y=|x|$,
              and one parabola that moved the ``wrong'' way & \blank{1.2cm} \\
2 & Guided Notes \& Practice & One curve in three costumes, $x^2-2x-3=(x-1)^2-4=(x+1)(x-3)$:
              what each form hands you for free. Four rows of notes and problems, then one new
              curve worked together & \blank{1.2cm} \\
3 & Homework & Reading vertex form, the inside-versus-outside pair, a graph to read and write
              as a rule, intercept form and its special case, a stone arch, and one
              multiple-choice item --- \textbf{scored; due the first class after two study
              halls} & \blank{1.2cm} \\
\midrule
  & \multicolumn{2}{r}{\textbf{Total}} & \blank{1.2cm} \\
\end{tabularx}
\end{tocbox}

\vspace{0.10in}

\boxguard[8]
\begin{remindbox}
\small
Every parabola is the \textbf{parent function} $y=x^2$ after a \textbf{transformation}, and one
curve can be written three ways --- \emph{each form hands you a different feature for free}.
\textbf{Vertex form} $a(x-h)^2+k$: vertex $(h,k)$, axis $x=h$. \textbf{Standard form}
$ax^2+bx+c$: $y$-intercept $(0,c)$, axis $x=-\frac{b}{2a}$, then substitute back for the vertex.
\textbf{Intercept form} $a(x-p)(x-q)$: \textbf{zeros} $p$ and $q$, axis at their midpoint. In
every form $a$ gives direction and width, and $a<0$ makes the vertex a \emph{maximum}.
\textbf{Mind the sign inside the bracket:} $y=(x-1)^2$ moves the parent \emph{right} $1$ while
$y=x^2-1$ moves it \emph{down} $1$ --- but as a \emph{factor}, $(x+1)$ means the zero $x=-1$.
\end{remindbox}

\end{document}
